

% =============================================================================
% Thermodynamics of goethite dehydroxylation, consolidated at the end of the
% Discussion (moved 2026-09-29): model, stability, and nucleation control.
% =============================================================================
% --- Added 2026-09-25 (goethite stability integration; see CHANGES_JOURNAL.md) ---
% Figures for the goethite stability/kinetics paragraphs below and in the
% Discussion are produced by Chapter2/Figures/fig_goethite_hematite_phase_diagram.py,
% fig_goethite_vacuum_kinetics.py and fig_goethite_hydrothermal_tests.py.
\paragraph{Thermodynamic Model of Goethite Dehydroxylation}\label{sec:gt_thermo_methods}

The preceding sections establish from the release patterns alone that the vacuum step heats of
Win-06-01b record conversion to hematite, whereas the hydrothermal step heat records diffusion
without detectable dehydration. Here we use the thermodynamics of goethite dehydroxylation to
explain why water has this effect (\S\ref{sec:gt_thermo}) and to estimate how much solid-state
conversion could have occurred in the hydrothermal experiment (\S\ref{sec:gt_nucleation}).

Goethite breaks down to hematite by the dehydroxylation reaction
\begin{equation}
  2\,\alpha\text{-FeOOH} \;\rightleftharpoons\; \alpha\text{-Fe}_2\text{O}_3 + \text{H}_2\text{O}_{(\text{fluid})},
  \label{eq:gt_rxn}
\end{equation}
whose Gibbs energy change, written per mole of hematite formed, is
\begin{equation}
  \Delta G_r(T,P) = G_{\mathrm{hem}}(T,P) + G_{\mathrm{w}}(T,P) - 2\,G_{\mathrm{goe}}(T,P) + \Delta G_{\mathrm{surf}} .
  \label{eq:gt_dGr}
\end{equation}
Goethite is stable where $\Delta G_r>0$ and hematite plus water where $\Delta G_r<0$.

\textit{Gibbs energies of the solids.} The Gibbs energy of a phase at $(T,P)$ follows from
integrating $dG = -S\,dT + V\,dP$ from the reference state $(T_r,P_r) = (298.15~\mathrm{K}, 1~\mathrm{bar})$,
first in temperature at $P_r$ and then in pressure at $T$. At $P_r$,
$H(T) = H(T_r) + \int_{T_r}^{T}C_p\,dT$ and $S(T) = S(T_r) + \int_{T_r}^{T}(C_p/T)\,dT$, so
$G(T,P_r) = H(T) - T S(T)$. Adding $\int_{P_r}^{P}V\,dP$, and treating the solids as incompressible
with no thermal expansion, gives
\begin{equation}
  G_i(T,P) = \Delta_f H_i^{0} + \int_{T_r}^{T} C_{p,i}\,dT
  \;-\; T\left[S_i^{0} + \int_{T_r}^{T}\frac{C_{p,i}}{T}\,dT\right]
  \;+\; V_i\,(P-P_r),
  \label{eq:gt_Gsolid}
\end{equation}
where we have set $H(T_r) = \Delta_f H_i^{0}$, the enthalpy of formation from the elements, and
$S(T_r) = S_i^{0}$, the third-law entropy. This is the ``apparent'' Gibbs energy of formation
\citep{berman1988}; see also \citet{helgeson1978} and \citet{anderson2005}. It differs from
the true Gibbs energy of formation,
$\Delta_f G_i = \Delta_f H_i - T(S_i - \sum_e n_{e,i} S_e)$, by the term $T\sum_e n_{e,i}S_e$,
where $n_{e,i}$ is the number of atoms of element $e$ in phase $i$ and $S_e$ the entropy of the
element in its reference state. For any balanced reaction $\sum_i \nu_i n_{e,i} = 0$ for every element,
so these terms cancel in $\Delta G_r$. The incompressible approximation affects $\Delta G_r$
through $\Delta V_s (P-P_r)$ with $\Delta V_s = V_{\mathrm{hem}} - 2V_{\mathrm{goe}} = -11.5$~cm$^3$,
i.e.\ at most a few kJ below 1~GPa; compressibility and thermal expansion change $V_i$ by
$\lesssim1\%$ over the conditions considered and are neglected.

Goethite heat capacity uses the Berman--Brown form \citep{berman1985} refined by
\citet{majzlan2003},
\begin{equation}
  C_p = k_0 + k_1 T^{-1/2} + k_2 T^{-2} + k_3 T^{-3},
\end{equation}
which integrates to
\begin{align}
  \int_{T_r}^{T} C_p\,dT &= k_0 (T-T_r) + 2k_1\left(T^{1/2}-T_r^{1/2}\right) \notag\\
    &\quad - k_2\left(T^{-1}-T_r^{-1}\right) - \tfrac{1}{2}k_3\left(T^{-2}-T_r^{-2}\right), \\
  \int_{T_r}^{T} \frac{C_p}{T}\,dT &= k_0 \ln\frac{T}{T_r} - 2k_1\left(T^{-1/2}-T_r^{-1/2}\right) \notag\\
    &\quad - \tfrac{1}{2}k_2\left(T^{-2}-T_r^{-2}\right) - \tfrac{1}{3}k_3\left(T^{-3}-T_r^{-3}\right).
\end{align}
Hematite uses a Maier--Kelley form, $C_p = a + bT + cT^{-2}$, with
\begin{align}
  \int_{T_r}^{T} C_p\,dT &= a(T-T_r) + \tfrac{1}{2}b\left(T^2-T_r^2\right) - c\left(T^{-1}-T_r^{-1}\right), \\
  \int_{T_r}^{T} \frac{C_p}{T}\,dT &= a\ln\frac{T}{T_r} + b(T-T_r) - \tfrac{1}{2}c\left(T^{-2}-T_r^{-2}\right).
\end{align}
% TODO: cite the source of the hematite Maier-Kelley coefficients (Kelley-type fit,
% cf. Robie & Hemingway 1995, USGS Bull. 2131); Majzlan et al. used Berman (1988). Differences are < 1 J/K/mol at 298 K.

\textit{Water.} Water is evaluated with the IAPWS-95 equation of state \citep{wagner2002},
implemented in CoolProp \citep{bell2014}, which is valid for liquid, vapor and supercritical
water to 1~GPa. In the same apparent convention,
\begin{equation}
  G_{\mathrm{w}}(T,P) = \Delta_f H_{\mathrm{w}}^{0} + \left[h(T,P) - h(T_r,P_r)\right]
  - T\left[S_{\mathrm{w}}^{0} + s(T,P) - s(T_r,P_r)\right],
  \label{eq:gt_Gw}
\end{equation}
where $h$ and $s$ are the IAPWS-95 molar enthalpy and entropy and the reference state is liquid
water at $T_r,P_r$. Equation~\eqref{eq:gt_Gw} is exact within the equation of state; no ideal-gas
approximation is made anywhere in the calculation. Table~\ref{tab:gt_thermo_data} lists all
input data. We use the enthalpies and entropies of \citet{majzlan2003} refined by mathematical
programming against the phase-equilibrium brackets of \citet{voigt1981}.

\begin{table}[htbp]
  \centering
  \caption{Thermodynamic data used for Eq.~\eqref{eq:gt_dGr}.}
  \label{tab:gt_thermo_data}
  \begin{footnotesize}
  \begin{tabular}{lcc}
    \hline
    & Goethite ($\alpha$-FeOOH) & Hematite ($\alpha$-Fe$_2$O$_3$) \\
    \hline
    $\Delta_f H^0$ (kJ\,mol$^{-1}$) & $-561.4$ & $-825.8$ \\
    $S^0$ (J\,mol$^{-1}$\,K$^{-1}$) & 59.5 & 87.4 \\
    $V^0$ (cm$^3$\,mol$^{-1}$) & 20.881 & 30.274 \\
    $C_p$ coefficients & $k_0=133.8765$ & $a=98.28$ \\
          & $k_1=-738.5577$ & $b=7.782\times10^{-2}$ \\
          & $k_2=-2.177657\times10^{6}$ & $c=-1.485\times10^{6}$ \\
          & $k_3=2.165124\times10^{8}$ & \\
    $C_p(298\ \mathrm{K})$ (J\,mol$^{-1}$\,K$^{-1}$) & 74.8 & 104.8 \\
    $\gamma_{\mathrm{hydrated}}$ (J\,m$^{-2}$) & $0.60\pm0.10$ & $0.75\pm0.16$ \\
    $\gamma_{\mathrm{anhydrous}}$ (J\,m$^{-2}$) & $0.91\pm0.09$ & $1.90\pm0.30$ \\
    \hline
  \end{tabular}
  \end{footnotesize}

  \vspace{0.5ex}
  \begin{minipage}{0.85\linewidth}\footnotesize
  Enthalpy, entropy and $C_p$: \citet{majzlan2003a,majzlan2003}; surface energies: \citet{navrotsky2008}.
  Water: $\Delta_f H^0_{\mathrm{w}} = -285.83$~kJ\,mol$^{-1}$, $S^0_{\mathrm{w}} = 69.95$~J\,mol$^{-1}$\,K$^{-1}$; IAPWS-95 for $h(T,P)$, $s(T,P)$.
  \end{minipage}
\end{table}

As a check, at 298.15~K and 1~bar with liquid water, Eqs.~\eqref{eq:gt_dGr}--\eqref{eq:gt_Gw}
give $\Delta G_r = -0.26$~kJ: hematite plus liquid water is favored by only 0.13~kJ per mole of
FeOOH. This is well inside the $\pm1.4$~kJ\,mol$^{-1}$ uncertainty that \citet{majzlan2003}
assign to the goethite--hematite free-energy difference; the two assemblages are essentially
balanced at ambient conditions. We carry this uncertainty ($\pm2.8$~kJ per reaction as written)
as bounding curves on the phase diagram.

\textit{Surface energy.} A crystal of finite size carries an excess free energy at its
surface in addition to the free energy of its bulk \citep{gibbs1878}. For a population of
particles of fixed size, this excess is $\gamma A$, where $\gamma$ is the surface energy
(J\,m$^{-2}$) and $A$ the total surface area. Both the bulk and surface terms scale with the
amount of material at fixed particle size, so the surface contribution can be written per mole,
$G = G_{\mathrm{bulk}} + \gamma A_m$, with $A_m$ the surface area per mole. $A_m$ grows as $1/d$,
so the term is negligible for coarse crystals and significant for nanoparticles
\citep{navrotsky2008}. For the reaction as written (one mole of hematite from two moles of
goethite),
\begin{equation}
  \Delta G_{\mathrm{surf}} = \gamma_{\mathrm{hem}} A_{m,\mathrm{hem}} - 2\,\gamma_{\mathrm{goe}} A_{m,\mathrm{goe}},
  \label{eq:gt_dGsurf}
\end{equation}
where $A_{m,\mathrm{hem}}$ is per mole of Fe$_2$O$_3$ and $A_{m,\mathrm{goe}}$ per mole of FeOOH.

A size must be chosen for each phase. We assume pseudomorphic dehydroxylation, in which hematite
inherits the external envelope of the parent goethite, so that both phases share the same
external surface area per mole of Fe. This is also the basis of the stability diagrams of
\citet{navrotsky2008}. For spherical goethite grains of diameter $d$ this shared area is
\begin{equation}
  A_{\mathrm{Fe}} = \frac{6\,V_{\mathrm{goe}}}{d}\quad\text{per mole of Fe}.
\end{equation}
Goethite has one Fe per formula unit, so $A_{m,\mathrm{goe}} = A_{\mathrm{Fe}}$. Hematite has two,
so $A_{m,\mathrm{hem}} = 2A_{\mathrm{Fe}}$. Substituting into Eq.~\eqref{eq:gt_dGsurf},
\begin{equation}
  \Delta G_{\mathrm{surf}} = \gamma_{\mathrm{hem}}\,(2A_{\mathrm{Fe}}) - 2\,\gamma_{\mathrm{goe}}\,A_{\mathrm{Fe}}
  = 2A_{\mathrm{Fe}}\left(\gamma_{\mathrm{hem}} - \gamma_{\mathrm{goe}}\right).
  \label{eq:gt_equalarea}
\end{equation}
The factor of two counts the two moles of Fe that pass through the reaction. It multiplies both
terms because both phases carry the same area per Fe. For $d=10$~nm,
$A_{\mathrm{Fe}} = 6\times20.881\times10^{-6}/10^{-8} = 1.253\times10^{4}$~m$^2$\,mol$^{-1}$, and
with hydrated surface energies
$\Delta G_{\mathrm{surf}} = 2\times1.253\times10^{4}\times(0.75-0.60) = +3.76$~kJ.
The corresponding values are $+0.38$~kJ at 100~nm and $+0.04$~kJ at 1~\textmu m. A positive
$\Delta G_{\mathrm{surf}}$ stabilizes goethite. The result depends on the choice of basis: at
equal diameter the same hydrated surface energies destabilize fine goethite
(\S\ref{sec:gt_supp_basis}).

We use hydrated surface energies. They are referenced to liquid water, so they apply wherever the
surfaces are in equilibrium with water at or near saturation, which includes all of the
hydrothermal experiments. Anhydrous surface energies apply only to surfaces that have lost their
chemisorbed water, i.e.\ under vacuum.

\textit{Phase diagram.} The phase diagram (Fig.~\ref{fig:gt_PT}) is computed by solving
$\Delta G_r(T,P) = 0$ (Eqs.~\eqref{eq:gt_dGr}--\eqref{eq:gt_equalarea}) numerically for
$P = P_{\mathrm{H_2O}}$ at each temperature, using Brent's method on $\log_{10}P$ between
$10^{-12}$ and $10^{4}$~bar, with the solids at the fluid pressure. Where no root exists below the
melting curve of ice, the curve is truncated.

\textit{Driving force and the role of water.} The quantity that matters for kinetics is the
driving force for hematite formation, $\Delta\mu \equiv -\Delta G_r^{0}$ (per mole of Fe$_2$O$_3$).
Water controls $\Delta\mu$ through its own chemical potential, because it is a product of
reaction~\eqref{eq:gt_rxn}. $G_{\mathrm{w}}$ in Eq.~\eqref{eq:gt_dGr} is the chemical potential of
H$_2$O at the prevailing water fugacity. It depends on the partial pressure of water, not on the
total pressure. An inert gas, or any other mechanical load, acts only through the solid-volume
term $\Delta V_s(P-P_r)$, which contributes $-1.15$~J\,bar$^{-1}$ and slightly favors the denser
hematite. At 40~bar that is 46~J, three orders of magnitude smaller than the effect of 40~bar of
water. The water-buffered step-heating technique therefore stabilizes goethite primarily by
mass action on the bulk reaction. Surface hydration is a second-order correction that matters
only for nanometer grains (Eq.~\eqref{eq:gt_equalarea}).

% --- Added 2026-09-25 (goethite stability integration; see CHANGES_JOURNAL.md) ---
\paragraph{Thermodynamic Stability of Goethite under Experimental Conditions}\label{sec:gt_thermo}

The thermodynamic model of \S\ref{sec:gt_thermo_methods} places the goethite--hematite--water
boundary in $\log P_{\mathrm{H_2O}}$--$T$ space (Fig.~\ref{fig:gt_PT}). Above the
saturation pressure, water is liquid, the reaction volume
$\Delta V_r = V_{\mathrm{hem}} + V_{\mathrm{w}} - 2V_{\mathrm{goe}} \approx +6.6$~cm$^3$ is small,
and the boundary becomes nearly vertical in $\log P$--$T$ space. Bulk goethite meets the
liquid--vapor curve at 18~$^\circ$C, 100~nm goethite at 28~$^\circ$C and 10~nm goethite at
108~$^\circ$C. At 25~$^\circ$C, bulk goethite requires $P_{\mathrm{H_2O}}\approx10^{2.6}$~bar
(compressed liquid). By contrast, 100~nm goethite is stable above $10^{-1.52}$~bar (96\%
relative humidity) and 10~nm goethite above $10^{-2.11}$~bar (24\% relative humidity). Bulk
goethite is therefore metastable at the conditions of every experiment considered here, in
water or in vacuum.

\begin{figure}[htbp]
  \centering
  \includegraphics[width=0.7\linewidth]{goethite_hematite_Ph2o_T_hydrated.pdf}
  \caption{Goethite--hematite--water equilibrium in $\log P_{\mathrm{H_2O}}$--$T$ space for bulk
  goethite and for 100 and 10~nm grains with hydrated surfaces, compared at equal surface area
  per mole of Fe (Eq.~\eqref{eq:gt_equalarea}). Goethite is stable above and to the left of each
  curve. Dash-dot curves: bulk boundary shifted by the $\pm1.4$~kJ\,mol$^{-1}$ FeOOH uncertainty of
  \citet{majzlan2003}. Dotted: liquid--vapor curve of water. Thermodynamic data in
  Table~\ref{tab:gt_thermo_data}.}
  \label{fig:gt_PT}
\end{figure}

Stability alone does not determine what happens in a step heat; the rate of conversion does,
and that depends on the driving force $\Delta\mu$. Table~\ref{tab:gt_dmu} gives the bulk driving
forces at the experimental temperatures, evaluated with IAPWS-95 water. At saturation,
$\Delta\mu$ is 6--11~kJ, only 1.6--2.4 times $RT$. Under the vacuum of a step-heating line
($P_{\mathrm{H_2O}}\approx10^{-8}$~atm) it is 76--106~kJ, about ten times larger. The chemical
potential of water is therefore the key variable separating the two experimental settings.

\begin{table}[htbp]
  \centering
  \caption{Bulk driving force for hematite formation, $\Delta\mu = -\Delta G_r$ (kJ per mole Fe$_2$O$_3$), at water saturation and at $P_{\mathrm{H_2O}} = 10^{-8}$~atm.}
  \label{tab:gt_dmu}
  \begin{tabular}{cccccc}
    \hline
    $T$ ($^\circ$C) & $P_{\mathrm{sat}}$ (bar) & $\Delta\mu_{\mathrm{sat}}$ & $\Delta\mu_{\mathrm{sat}}/RT$ & $\Delta\mu_{\mathrm{vac}}$ & $(\Delta\mu_{\mathrm{vac}}/\Delta\mu_{\mathrm{sat}})^2$ \\
    \hline
    150 & 4.76  & 5.68  & 1.61 & 75.8  & 178 \\
    175 & 8.93  & 6.88  & 1.84 & 83.4  & 147 \\
    200 & 15.55 & 8.10  & 2.06 & 91.0  & 126 \\
    225 & 25.50 & 9.36  & 2.26 & 98.5  & 111 \\
    250 & 39.76 & 10.65 & 2.45 & 106.0 & 99  \\
    \hline
  \end{tabular}
\end{table}

\paragraph{Water Suppresses Solid-State Conversion of Goethite}\label{sec:gt_nucleation}

\textit{Nucleation barrier.} In classical nucleation theory \citep{gibbs1878,volmer1926,becker1935},
reviewed by \citet{christian2002}, \citet{kashchiev2000} and, for mineral systems,
\citet{deyoreo2003} and \citet{lasaga1998}, a spherical hematite nucleus of radius $r$ has free
energy
\begin{equation}
  \Delta G(r) = -\frac{4}{3}\pi r^{3}\,\frac{\Delta\mu}{V_{\mathrm{hem}}} + 4\pi r^{2}\gamma,
\end{equation}
where $\Delta\mu$ is the driving force per mole of Fe$_2$O$_3$ (Table~\ref{tab:gt_dmu}),
$V_{\mathrm{hem}}$ the molar volume and $\gamma$ the energy of the interface surrounding the
nucleus. The first term is the free energy released by transforming the nucleus volume. The second
is the cost of creating its interface. Setting $d\Delta G/dr=0$ gives the barrier height
\begin{equation}
  \Delta G^{*} = \frac{16\pi\,\gamma^{3}\,V_{\mathrm{hem}}^{2}}{3\,\Delta\mu^{2}} .
  \label{eq:gt_barrier}
\end{equation}
Written per formula unit, $V_{\mathrm{hem}}\to V_{\mathrm{hem}}/N_A$ and $\Delta\mu\to\Delta\mu/N_A$, so
Avogadro's number cancels and Eq.~\eqref{eq:gt_barrier} is directly the barrier per nucleus in
joules. In the first-order limit of \S\ref{sec:gt_vacuum}, the rate constant is a nucleation
rate,
\begin{gather}
  k = A\exp\!\left(-\frac{E_a}{RT}\right)\exp\!\left(-\frac{\Delta G^{*}}{k_B T}\right), \notag\\
  \frac{\Delta G^{*}}{k_B T} = \frac{C\gamma^{3}}{T\,\Delta\mu^{2}},
  \qquad
  C \equiv \frac{16\pi V_{\mathrm{hem}}^{2}}{3k_B} = 1.112\times10^{15}~\mathrm{m^6\,K\,J^{-1}\,mol^{-2}} ,
  \label{eq:gt_knuc}
\end{gather}
where $A\exp(-E_a/RT)$ describes the atomic mobility and $\exp(-\Delta G^{*}/k_BT)$ the probability of
forming a critical nucleus. The barrier scales as $\gamma^3/\Delta\mu^2$. Between vacuum and water
saturation, $\Delta\mu$ falls tenfold and the barrier rises by the factor in the last column of
Table~\ref{tab:gt_dmu}: about 100 at 250~$^\circ$C and nearly 180 at 150~$^\circ$C.

\textit{An upper bound on the vacuum barrier.} The vacuum conversion rate is observed, so the
barrier that controls it cannot exceed what the fastest conceivable prefactor allows. In
transition-state theory \citep{eyring1935} no elementary step is faster than $k_BT/h$, which is
$1.09\times10^{13}$~s$^{-1}$ at 250~$^\circ$C. If every one of the $N$ formula units in a grain were a
nucleation site attempting at this frequency, the per-grain rate could not exceed
$N(k_BT/h)\exp(-\Delta G^{*}/k_BT)$. A 1~\textmu m goethite sphere contains
$N = (\pi d^3/6)/(V_{\mathrm{goe}}/N_A) = 1.5\times10^{10}$ formula units. Requiring this ceiling to
reach the observed $k_{\mathrm{vac}} = \ChTwoKinKTwoFifty$~s$^{-1}$ at 250~$^\circ$C (Eq.~\eqref{eq:gt_line})
gives
\begin{equation}
  \frac{\Delta G^{*}_{\mathrm{vac}}}{k_BT} \le \ln\frac{N\,k_BT/h}{k_{\mathrm{vac}}}
  = \ln\frac{1.5\times10^{10}\times1.09\times10^{13}}{\ChTwoKinKTwoFifty} = \ChTwoKinCeiling .
  \label{eq:gt_ceiling}
\end{equation}
A tenfold larger grain raises the bound by only $\ln 1000 = 7$. Real prefactors are smaller, so the
true barrier is lower still.

\textit{Nucleation must be internal.} The surface energies of \citet{navrotsky2008} fail this
bound by a wide margin. If the nucleus were bounded by a free, anhydrous hematite surface
($\gamma = 1.90$~J\,m$^{-2}$, $\Delta\mu = 106.0$~kJ at 250~$^\circ$C), Eq.~\eqref{eq:gt_knuc} would give
$\Delta G^{*}/k_BT = 1.3\times10^{3}$. Nucleation on a goethite surface lowers this. For a nucleus
shaped as a spherical cap on a flat substrate, the barrier is reduced by the purely geometric
factor
$f(\theta) = (2+\cos\theta)(1-\cos\theta)^2/4 \le 1$, where the contact angle $\theta$ is set by the
balance of interfacial tensions at the triple junction (Young's equation),
$\cos\theta = (\gamma_{\mathrm{goe}} - \gamma_{\mathrm{hem|goe}})/\gamma_{\mathrm{hem}}$
\citep{volmer1926,turnbull1950}. Even with a hematite--goethite interfacial energy of zero,
$\cos\theta \le 0.91/1.90 = 0.48$, $f\ge0.17$ and $\Delta G^{*}/k_BT\ge2.2\times10^{2}$. No
non-negative interfacial energy brings nucleation on free surfaces within the bound of
Eq.~\eqref{eq:gt_ceiling}. Hematite must therefore nucleate internally, where the only new
interface is between hematite and goethite.

This agrees with the known mechanism. Goethite dehydroxylates topotactically: the hematite
inherits the oxygen framework and orientation of the parent, and the pseudomorph develops pores
and twinned hematite domains \citep{limadefaria1963,watari1979,watari1983,cudennec2005}. Interfaces
of this kind are coherent or semicoherent. In the general classification of \citet{porter2009},
developed for metals, coherent interfaces have energies up to $\sim$0.2~J\,m$^{-2}$ and semicoherent
interfaces 0.2--0.5~J\,m$^{-2}$, compared with 0.5--1~J\,m$^{-2}$ for incoherent ones. For this system
specifically, density-functional calculations give 0.24~J\,m$^{-2}$ for the prismatic twin boundaries
between hematite domains formed by dehydrating goethite \citep{saito2016}. We are not aware of a
measured or calculated energy for the goethite--hematite interface itself. In what follows,
$\gamma$ is replaced by an effective interfacial energy $\gamma_{\mathrm{eff}}$, calibrated from the
data. It absorbs the geometry of internal nucleation and is a property of the sample.

\textit{Calibrating $\gamma_{\mathrm{eff}}$ from the trap-isolated steps.} Isolating the trap
raised $P_{\mathrm{H_2O}}$ in the line and lowered $\Delta\mu$, and the isothermal rate fell by a factor
$s$ (\S\ref{sec:gt_vacuum}). With $A$ and $E_a$ unchanged, Eq.~\eqref{eq:gt_knuc} gives
\begin{equation}
  \ln\frac{1}{s} = \frac{C\gamma_{\mathrm{eff}}^{3}}{T}\left(\frac{1}{\Delta\mu_{\mathrm{off}}^{2}} - \frac{1}{\Delta\mu_{\mathrm{on}}^{2}}\right)
  \quad\Longrightarrow\quad
  \gamma_{\mathrm{eff}} = \left[\frac{T\ln(1/s)}{C\left(\Delta\mu_{\mathrm{off}}^{-2} - \Delta\mu_{\mathrm{on}}^{-2}\right)}\right]^{1/3}.
  \label{eq:gt_gamma_cal}
\end{equation}
The driving forces are evaluated as follows.
\begin{enumerate}
  \item $\Delta\mu_{\mathrm{on}}$ and $\Delta\mu_{\mathrm{off}}$ are bulk driving forces,
  $-\Delta G_r$ from Eqs.~\eqref{eq:gt_dGr}--\eqref{eq:gt_Gw} with $\Delta G_{\mathrm{surf}} = 0$, at the
  step temperature and with the solids at vacuum.
  \item They differ only through the chemical potential of water: $\Delta\mu_{\mathrm{on}}$ is
  evaluated at $P_{\mathrm{H_2O}} = 10^{-8}$~atm and $\Delta\mu_{\mathrm{off}}$ at the water pressure
  reached with the trap isolated (below). At these pressures (at most 20~Pa), IAPWS-95 and the
  ideal gas agree, and
  $\Delta\mu_{\mathrm{on}} - \Delta\mu_{\mathrm{off}} = RT\ln(P_{\mathrm{off}}/P_{\mathrm{on}})$, which is
  32--41~kJ for the pressures in Table~\ref{tab:gt_gamma}.
  \item Surface terms are omitted in both. The barrier of an internal nucleus depends on the free
  energy per volume of material transformed, not on the external surface of the host grain. The
  grain or domain size of the natural sample is also not characterized. The external surface
  energy of the grain enters the overall thermodynamic stability (\S\ref{sec:gt_thermo}),
  but we do not include it in the barrier. The sensitivity to this choice is small
  (\S\ref{sec:gt_supp_barrier}).
\end{enumerate}
We bound the water pressure with the trap isolated by assuming that all evolved water was in the
vapor, distributed uniformly through an apparatus of volume $V_{\mathrm{app}}$ at 298~K:
\begin{equation}
  P_{\mathrm{H_2O}} = \frac{(m_{\mathrm{w}}/M_{\mathrm{w}})\,RT}{V_{\mathrm{app}}} .
\end{equation}
With $m_{\mathrm{w}} = 0.027$~mg at 230~$^\circ$C and 0.016~mg at 240~$^\circ$C ($1.50$ and
$0.89~\mu$mol), and $190\le V_{\mathrm{app}}\le600$~mL, Table~\ref{tab:gt_gamma} gives the pressures
and the resulting $\gamma_{\mathrm{eff}}$. The pressures are end-of-step values; we also evaluate half
of each as a time average for steady accumulation. Pressure enters only through $\ln P$ in
$\Delta\mu$ and then through a cube root, so a threefold uncertainty in volume changes
$\gamma_{\mathrm{eff}}$ by only about 10\%. The two steps agree within 15\%, and we adopt
\begin{equation}
  \gamma_{\mathrm{eff}} = 0.16\pm0.03~\mathrm{J\,m^{-2}} ,
\end{equation}
within the range expected for coherent and semicoherent interfaces. The corresponding barriers at
250~$^\circ$C are 0.5--1.3~$k_BT$ under vacuum and 52--129~$k_BT$ at saturation. Under vacuum the
barrier is negligible and the Arrhenius term of Eq.~\eqref{eq:gt_line} sets the rate; at saturation
the barrier dominates.

\begin{table}[htbp]
  \centering
  \caption{Water pressure with the trap isolated and the calibrated $\gamma_{\mathrm{eff}}$ (J\,m$^{-2}$; Eq.~\eqref{eq:gt_gamma_cal}).}
  \label{tab:gt_gamma}
  \begin{tabular}{cccccc}
    \hline
    $T$ ($^\circ$C) & $s$ & $V_{\mathrm{app}}$ (mL) & $P_{\mathrm{H_2O}}$ (Pa) & $\gamma_{\mathrm{eff}}$, end-of-step $P$ & $\gamma_{\mathrm{eff}}$, $P/2$ \\
    \hline
    230 & 0.333 & 190 & 19.6 & 0.138 & 0.145 \\
    230 & 0.333 & 600 & 6.2  & 0.150 & 0.158 \\
    240 & 0.224 & 190 & 11.6 & 0.164 & 0.172 \\
    240 & 0.224 & 600 & 3.7  & 0.178 & 0.188 \\
    \hline
  \end{tabular}
\end{table}

\textit{Translation to water saturation.} For the same sample, with the same internal interface
and the same mobility term, only $\Delta\mu$ changes between vacuum and saturation:
\begin{equation}
  \frac{k_{\mathrm{sat}}}{k_{\mathrm{vac}}} = \exp\!\left[-\frac{C\gamma_{\mathrm{eff}}^{3}}{T}\left(\frac{1}{\Delta\mu_{\mathrm{sat}}^{2}} - \frac{1}{\Delta\mu_{\mathrm{vac}}^{2}}\right)\right],
  \label{eq:gt_translate}
\end{equation}
with both driving forces evaluated as in the calibration (bulk, IAPWS-95 water at $P_{\mathrm{sat}}$ and
at $10^{-8}$~atm). For $\gamma_{\mathrm{eff}} = 0.14$, 0.16 and 0.19~J\,m$^{-2}$, this ratio is
$7\times10^{-23}$, $9\times10^{-34}$ and $5\times10^{-56}$ at 250~$^\circ$C, and $5\times10^{-43}$,
$7\times10^{-64}$ and $2\times10^{-106}$ at 200~$^\circ$C.

How large a barrier would still allow detectable conversion follows from the same kinetics. Over
the 442.8~h ($1.59\times10^{6}$~s) hydrothermal experiment, converting even 1\% of the sample at
250~$^\circ$C requires $k_{\mathrm{sat}}t\gtrsim0.01$, i.e.
\begin{equation}
  \frac{\Delta G^{*}_{\mathrm{sat}}}{k_BT} \lesssim \ln\frac{k_{\mathrm{vac}}\,t}{0.01}
  = \ln\frac{\ChTwoKinKTwoFifty\times1.59\times10^{6}}{0.01} = \ChTwoKinDetect ,
\end{equation}
far below the calibrated 52--129~$k_BT$. Applying Eq.~\eqref{eq:gt_translate} step by step to the
hydrothermal schedule, with $k_{\mathrm{vac}}$ from Eq.~\eqref{eq:gt_line}, the predicted conversion is
\begin{equation}
  \alpha \approx \sum_i k_{\mathrm{vac}}(T_i)\,\frac{k_{\mathrm{sat}}}{k_{\mathrm{vac}}}(T_i)\,\Delta t_i
  = 10^{\ChTwoKinLogAlphaLow},\ 10^{\ChTwoKinLogAlphaMid},\ 10^{\ChTwoKinLogAlphaHigh}
\end{equation}
for the three values of $\gamma_{\mathrm{eff}}$. Even the most permissive calibration,
$\gamma_{\mathrm{eff}} = 0.084$~J\,m$^{-2}$, gives $\alpha = \ChTwoKinAlphaPermissive$. That value would require
0.1~bar of water in the line, far above the volume constraint. If the observed release came from
conversion, it would imply $\alpha = F/(1-\phi) = \ChTwoKinHTFEnd/\ChTwoKinOneMinusPhi = \ChTwoKinHTAlphaObs$.

\textit{Assumptions.} This treatment is a consistent physical explanation of the data, not a
proof. It rests on four assumptions.
(i) The slowdown during the trap-isolated steps acts through $\Delta\mu$. If adsorbed water instead
slowed dehydroxylation by blocking reactive sites, these steps would not calibrate
$\gamma_{\mathrm{eff}}$.
(ii) The effective interface energy between hematite and goethite, $\gamma_{\mathrm{eff}}$, $A$ and $E_a$ are the same in vacuum and in water. The
mechanism is internal nucleation in both settings, so only the chemical potential of water
changes.
(iii) The barrier depends on the bulk driving force. Including the surface terms of
Eq.~\eqref{eq:gt_equalarea} for grain sizes down to 10~nm raises the predicted hydrothermal
conversion by about eight orders of magnitude, to $\alpha\approx10^{-23}$
(\S\ref{sec:gt_supp_barrier}).
(iv) Only solid-state conversion is modeled. Conversion by dissolution of goethite and
precipitation of hematite from the fluid is not described by Eq.~\eqref{eq:gt_translate}
(\S\ref{sec:gt_pigment}).
The comparison is valid only between experimental conditions on the same sample. The calibrated
$\gamma_{\mathrm{eff}}$ absorbs the sample's microstructure and should not be transferred to other
goethites.

Assumption (iv) matters for the synthetic pigment of \S\ref{sec:gt_pigment}. Grain size and bulk
thermodynamics do not explain why the pigment blades reacted when the natural goethite did not. At 1~\textmu m, the surface contribution to $\Delta G_r$ is only
$\sim$0.04~kJ, and the bulk driving force at saturation is the same 6--11~kJ for both materials
(Table~\ref{tab:gt_dmu}). The most likely difference is the fluid itself. Commercial iron-oxide
pigments are commonly produced by precipitation from iron-salt solutions. They can carry
residual soluble salts, such as sulfate, and may be surface-treated. Such contaminants would
change the pH and ionic strength of the water they are heated in, and with them the solubility
of Fe(III). The dissolution--precipitation pathway is limited by how much iron the fluid can
carry, so a modest increase in Fe(III) solubility could turn a negligible reaction into a
visible one. Differences in residue loading or in the handling of individual ampoules would also
explain why the outcome was not reproducible. The natural goethite, by contrast, was heated in
distilled water, where near-neutral Fe(III) solubility is very low.


\clearpage

% SUPPLEMENT --------------------------------------------------------


% =============================================================================
% Goethite dehydroxylation kinetics and stability: supplementary sections.
% Added 2026-09-25 (goethite stability integration; see CHANGES_JOURNAL.md).
% =============================================================================

% -----------------------------------------------------------------------------
\section{Kinetic mechanisms and the limits of a single heating schedule}
\label{sec:gt_supp_mech}
% -----------------------------------------------------------------------------

Solid-state rate laws take the form $d\alpha/dt = A\exp(-E_a/RT)\,h(\alpha)$. For any heating
schedule this integrates to
\begin{equation}
  g(\alpha) \equiv \int_0^{\alpha}\frac{d\alpha'}{h(\alpha')} = x,
  \qquad
  x = A\sum_j \Delta t_j\exp\!\left(-\frac{E_a}{RT_j}\right).
\end{equation}
We tested the mechanisms in Table~\ref{tab:gt_mech_forms}, including spherical diffusion on the
same footing. Under vacuum, the driving-force factor $1-\exp(-\Delta\mu/RT)$ is unity to better
than $10^{-9}$, so $x$ is simply an Arrhenius-weighted time. Diffusion and conversion then
differ only in the shape $\alpha(x)$ and in the fitted $E_a$.

\begin{table}[htbp]
  \centering
  \caption{Kinetic models and their integrated forms.}
  \label{tab:gt_mech_forms}
  \footnotesize
  \setlength{\tabcolsep}{4pt}
  \resizebox{\textwidth}{!}{%
  \begin{tabular}{llll}
    \hline
    Model & $h(\alpha)$ & $g(\alpha)$ & $\alpha(x)$ \\
    \hline
    R2 contracting cylinder & $2(1-\alpha)^{1/2}$ & $1-(1-\alpha)^{1/2}$ & $1-(1-x)^2$ \\
    R3 contracting sphere & $3(1-\alpha)^{2/3}$ & $1-(1-\alpha)^{1/3}$ & $1-(1-x)^3$ \\
    F1 first order & $1-\alpha$ & $-\ln(1-\alpha)$ & $1-e^{-x}$ \\
    A2, A3 Avrami ($n=2,3$) & $n(1-\alpha)[-\ln(1-\alpha)]^{1-1/n}$ & $[-\ln(1-\alpha)]^{1/n}$ & $1-e^{-x^n}$ \\
    D3 sphere diffusion & --- & --- & $1-\dfrac{6}{\pi^2}\displaystyle\sum_n \dfrac{e^{-n^2\pi^2 x}}{n^2}$ \\
    \hline
  \end{tabular}}
\end{table}

\paragraph{Synthetic test.} We generated noiseless release curves from the R3 model
($E_a = 150$~kJ\,mol$^{-1}$) on two schedules: a monotonic series of 1~h steps from 150 to
390~$^\circ$C, and a cycled schedule with retrograde steps and long isothermal holds. We then
fitted every model with $E_a$ and $A$ free. Table~\ref{tab:gt_synth} shows that every model
except the Avrami forms fits within 1\% in $F$ by adjusting $E_a$. With realistic analytical
uncertainty, goodness of fit to a single schedule cannot identify the mechanism. The fitted
activation energy trades off against the shape of $\alpha(x)$. Discrimination requires a second,
independent schedule, an independent measure of $\alpha$, or both. The vacuum experiments supply
both (\S\ref{sec:gt_vacuum}).

\begin{table}[htbp]
  \centering
  \caption{Fits to synthetic R3 data ($E_a$ in kJ\,mol$^{-1}$; rms misfit in $F$).}
  \label{tab:gt_synth}
  \begin{tabular}{lcccc}
    \hline
    & \multicolumn{2}{c}{Monotonic schedule} & \multicolumn{2}{c}{Cycled schedule} \\
    Model & $E_a$ & rms & $E_a$ & rms \\
    \hline
\ChTwoKinSynthRows
    \hline
  \end{tabular}
\end{table}

\paragraph{Fits to the vacuum data.} Table~\ref{tab:gt_mech_fits} gives the fit of each model to
the cumulative KAF $^3$He release within the fitting window. It also gives the misfit of the
resulting out-of-sample prediction of HDM, and the fit to the KAF water data. F1 gives the best
fit and the best prediction. R3 and R2 are close behind, so we describe the mechanism as
nucleation- or interface-controlled rather than claim first order specifically. The Avrami
fits reach the lower bound of $E_a$ and fail on both data sets. Spherical diffusion with
$\phi=0$ fits KAF \ChTwoKinDThreeRmsRatio{} times worse than F1 and cannot produce the $F\approx0.89$
plateau. To compare diffusion on the same footing, we also fit D3 with $\phi$ fixed at the F1 value
and with $\phi$ free, over the same \ChTwoKinNKAF{} KAF steps. The imposed plateau improves the KAF
misfit to \ChTwoKinDThreePhiRmsKAF{} and \ChTwoKinDThreeFreeRmsKAF, which still ranks below F1, R3 and
R2, and the out-of-sample prediction of HDM does not improve (rms \ChTwoKinDThreePhiRmsHDM{} and
\ChTwoKinDThreeFreeRmsHDM, against \ChTwoKinFOneRmsHDM{} for F1). These fits also require
$E_a$ of \ChTwoKinDThreePhiEa--\ChTwoKinDThreeFreeEa~kJ\,mol$^{-1}$ and a rate factor
$s = \ChTwoKinDThreePhiS$ for the trap-isolated steps, although a diffusive process should not
respond to the water pressure at all.

The fitting window excludes the low-temperature steps, where release is dominated by diffusion from
goethite, and the plateau, where it is dominated by diffusion from hematite; the rms misfits are
computed over the same steps. Repeating the fits with every step included gives the same result:
F1 remains the best model ($E_a = \ChTwoKinAllFOneEa$~kJ\,mol$^{-1}$, $\phi = \ChTwoKinAllPhi$, rms \ChTwoKinAllFOneRmsKAF{} for KAF and \ChTwoKinAllFOneRmsHDM{}
for the HDM prediction), followed by R3 (\ChTwoKinAllRThreeEa~kJ\,mol$^{-1}$; \ChTwoKinAllRThreeRmsKAF, \ChTwoKinAllRThreeRmsHDM) and R2 (\ChTwoKinAllRTwoEa~kJ\,mol$^{-1}$;
\ChTwoKinAllRTwoRmsKAF, \ChTwoKinAllRTwoRmsHDM), with D3 at \ChTwoKinAllDThreeEa~kJ\,mol$^{-1}$ (\ChTwoKinAllDThreeRmsKAF, \ChTwoKinAllDThreeRmsHDM;
\ChTwoKinAllDThreePhiRmsKAF{} and \ChTwoKinAllDThreePhiRmsHDM{} with $\phi$ fixed at the F1 value).

\begin{table}[htbp]
  \centering
  \caption{Kinetic and Diffusive Model fits to the vacuum experiments (from \citet{monteiro2017}~and~\citet{farley2024}). $s$: rate factor for the trap-isolated steps. All rms values are over the same \ChTwoKinNKAF{} KAF and \ChTwoKinNHDM{} HDM steps. D3 is fitted with $\phi$ fixed at 0, fixed at the F1 value, and free; the water fit does not depend on $\phi$ and is given once.}
  \label{tab:gt_mech_fits}
  \begin{tabular}{lccccccc}
    \hline
    Model & $E_a$ & $\ln A$ & $\phi$ & $s$ & rms KAF & rms HDM & $E_a$ (water) \\
    \hline
\ChTwoKinMechFitRows
    \hline
    \multicolumn{8}{l}{\footnotesize $E_a$ in kJ\,mol$^{-1}$, $A$ in s$^{-1}$. $^{a}$At the lower bound of the fit.}
  \end{tabular}
\end{table}


% -----------------------------------------------------------------------------
\section{Sensitivity of the size effect to the surface-area basis}
\label{sec:gt_supp_basis}
% -----------------------------------------------------------------------------

The sign of the surface stabilization depends on how the surface areas of the two phases are
related. At equal area per mole of Fe (Eq.~\eqref{eq:gt_equalarea}), goethite is stabilized
whenever $\gamma_{\mathrm{hem}}>\gamma_{\mathrm{goe}}$. If both phases are instead spheres of the
same diameter $d$, then $A_i = 6V_i/d$ and
\begin{equation}
  \Delta G_{\mathrm{surf}} = \frac{6}{d}\left(\gamma_{\mathrm{hem}}V_{\mathrm{hem}} - 2\gamma_{\mathrm{goe}}V_{\mathrm{goe}}\right),
\end{equation}
which stabilizes goethite only if
\begin{equation}
  \frac{\gamma_{\mathrm{hem}}}{\gamma_{\mathrm{goe}}} > \frac{2V_{\mathrm{goe}}}{V_{\mathrm{hem}}} = \frac{2\times20.881}{30.274} = 1.379 .
\end{equation}
Hematite is denser and carries less surface per Fe at equal size. For hydrated surfaces,
$\gamma_{\mathrm{hem}}/\gamma_{\mathrm{goe}} = 0.75/0.60 = 1.25 < 1.379$, so fine goethite is
destabilized: at 10~nm,
$\Delta G_{\mathrm{surf}} = (6/10^{-8})(0.75\times30.274 - 2\times0.60\times20.881)\times10^{-6} = -1.41$~kJ.
For anhydrous surfaces, $1.90/0.91 = 2.09 > 1.379$ and $\Delta G_{\mathrm{surf}} = +11.7$~kJ at
10~nm. Table~\ref{tab:gt_basis} compares the two bases. The equal-area basis corresponds to
pseudomorphic conversion and underlies the main-text calculations. The equal-diameter basis is
more appropriate where hematite grows as separate crystals, as in fluid-mediated conversion
(\S\ref{sec:gt_pigment}). There the surface term slightly favors hematite, but it is
negligible at micrometer sizes.

\begin{table}[htbp]
  \centering
  \caption{$\Delta G_{\mathrm{surf}}$ (kJ per reaction as written) for the two surface-area bases.}
  \label{tab:gt_basis}
  \begin{tabular}{lcccc}
    \hline
    & \multicolumn{2}{c}{Hydrated} & \multicolumn{2}{c}{Anhydrous} \\
    $d$ & equal area & equal diameter & equal area & equal diameter \\
    \hline
    1~$\mu$m & $+0.04$ & $-0.01$ & $+0.25$ & $+0.12$ \\
    100~nm & $+0.38$ & $-0.14$ & $+2.48$ & $+1.17$ \\
    10~nm & $+3.76$ & $-1.41$ & $+24.8$ & $+11.7$ \\
    \hline
  \end{tabular}
\end{table}

\begin{figure}[htbp]
  \centering
  \includegraphics[width=\linewidth]{size_basis_sensitivity}
  \caption{Goethite--hematite--water boundaries for hydrated surfaces at (a) equal surface area
  per mole of Fe and (b) equal particle diameter. The size effect changes sign between the two
  bases.}
  \label{fig:gt_basis}
\end{figure}


% -----------------------------------------------------------------------------
\section{Nucleation barrier ratio and the calibration of $\gamma_{\mathrm{eff}}$}
\label{sec:gt_supp_barrier}
% -----------------------------------------------------------------------------

From Eq.~\eqref{eq:gt_barrier}, the ratio of barriers at water saturation and under vacuum at
the same temperature is
\begin{equation}
  \frac{\Delta G^{*}_{\mathrm{sat}}}{\Delta G^{*}_{\mathrm{vac}}}
  = \left(\frac{\Delta\mu_{\mathrm{vac}}}{\Delta\mu_{\mathrm{sat}}}\right)^{2}
    \left(\frac{\gamma_{\mathrm{sat}}}{\gamma_{\mathrm{vac}}}\right)^{3}.
\end{equation}
The prefactors, the temperature and the molar volume cancel. The first factor follows from the
thermodynamic data alone (Table~\ref{tab:gt_dmu}; 99 at 250~$^\circ$C for bulk goethite). The
second is 1 if nucleation occurs at the same internal interface in both settings. It is
$(0.75/1.90)^3 = 1/16$ in the extreme case that the relevant interface is a free hematite
surface, hydrated at saturation and anhydrous under vacuum. Figure~\ref{fig:gt_barrier_ratio}
shows the ratio between these bounds for bulk and 10~nm goethite. Even at the lower bound, the
saturated barrier is about six times the vacuum barrier at 250~$^\circ$C. Because the barrier
appears in an exponent, that is still a difference of many orders of magnitude in rate.

\begin{figure}[htbp]
  \centering
  \includegraphics[width=0.7\linewidth]{nucleation_barrier_ratio_vs_T}
  \caption{Ratio of hematite nucleation barriers at water saturation and at
  $P_{\mathrm{H_2O}}=10^{-8}$~atm for bulk and 10~nm goethite. Solid: same interface in both
  settings; dashed: free hematite surfaces (hydrated vs anhydrous). Below the vertical dotted
  lines, goethite is stable at saturation and there is no driving force.}
  \label{fig:gt_barrier_ratio}
\end{figure}

Figure~\ref{fig:gt_translation}a shows $\gamma_{\mathrm{eff}}$ from Eq.~\eqref{eq:gt_gamma_cal} as
a function of the assumed water pressure during the trap-isolated steps. It spans
0.08--0.20~J\,m$^{-2}$ over $10^{-5}$--$10^{-1}$~bar and 0.14--0.19~J\,m$^{-2}$ over the
volume-constrained range (Table~\ref{tab:gt_gamma}). Figure~\ref{fig:gt_translation}b shows the
predicted slowdown at saturation (Eq.~\eqref{eq:gt_translate}) for the calibrated values and for
the most permissive value.

\begin{figure}[htbp]
  \centering
  \includegraphics[width=\linewidth]{hydrothermal_translation}
  \caption{(a) $\gamma_{\mathrm{eff}}$ calibrated from the two trap-isolated KAF steps as a
  function of the assumed $P_{\mathrm{H_2O}}$; grey band: range allowed by the apparatus volume
  (190--600~mL) and the water released in each step. (b) Ratio of conversion rate constants at
  water saturation and under vacuum (Eq.~\eqref{eq:gt_translate}); grey band: temperatures of the
  hydrothermal experiments.}
  \label{fig:gt_translation}
\end{figure}

\paragraph{Sensitivity to surface terms.} The calibration and translation use bulk driving
forces (\S\ref{sec:gt_nucleation}). To test this choice, we repeat both with the surface
term of Eq.~\eqref{eq:gt_equalarea} included in $\Delta\mu$. We use anhydrous surfaces for the vacuum
and trap-isolated steps and hydrated surfaces at saturation, the combination that most reduces the
vacuum driving force relative to the saturated one. Table~\ref{tab:gt_surf_sens} shows that the
conclusion is unchanged down to 10~nm grains.

\begin{table}[htbp]
  \centering
  \caption{Effect of including surface terms in $\Delta\mu$ on the calibration and the predicted hydrothermal conversion. Driving forces in kJ per mole Fe$_2$O$_3$ at 250~$^\circ$C; $\gamma_{\mathrm{eff}}$ range over Table~\ref{tab:gt_gamma}.}
  \label{tab:gt_surf_sens}
  \begin{tabular}{lcccc}
    \hline
    Grain size & $\Delta\mu_{\mathrm{vac}}$ & $\Delta\mu_{\mathrm{sat}}$ & $\gamma_{\mathrm{eff}}$ (J\,m$^{-2}$) & $\log_{10}\alpha_{\mathrm{HT}}$ \\
    \hline
    Bulk & 106.0 & 10.65 & 0.138--0.188 & $-32.0$ \\
    100~nm & 103.5 & 10.27 & 0.133--0.182 & $-31.2$ \\
    10~nm & 81.2 & 6.89 & 0.090--0.131 & $-23.3$ \\
    \hline
    \multicolumn{5}{l}{\footnotesize $\alpha_{\mathrm{HT}}$: predicted conversion over the hydrothermal schedule, median $\gamma_{\mathrm{eff}}$.}
  \end{tabular}
\end{table}



